\section{Données du Challenge RAG@EvalLLM}

    \begin{verbatim}
Présentation des données : 
- pdf à traiter sans structure normée du contenu
- majoritairement en français sauf quelques uns en anglais
- peuvent contenir des images
Présentation de la tâche du challenge (sauf si fait dans introduction)
- Identification des documents pertinents pour une question (phrase 
parfois complexe ou suite de mots clés) 
- questions fournies au dernier moment
    \end{verbatim}

\section{Structure de données visée}
    \begin{verbatim}
Base de données:        
- Paires 
- liste des champs et justification
Choix pour datasets
- Dataset avec des paires question-document (questions différentes et 
différents pdf, taille des documents, répartition positif/négatif)
- Dataset avec une seule question (construction à partir de 3 pdf) 
[description du dataset dans les expés et les specs]
    - motivation (reprendre l'expé présentée par MechIR, faciliter la création)
    - limites (biais, diversité ~> généralisabilité)
=> Placer le nom des outils utilisés (PyMuPDF4LLM et langchain.text_splitter.MarkdownSplitter)

    \end{verbatim}
